\documentclass[11pt]{article}
\usepackage{fontspec}
\usepackage{xeCJK}
\usepackage{booktabs}
\usepackage{hyperref}
\setmainfont{Libertinus Sans}
\title{上古漢語藏文轉寫方案}
\author{Kelvin Yang}
\date{2026年8月}
\begin{document}
\maketitle

本方案爲以藏文拼寫上古漢語之方案，以 Baxter--Sagart 擬音爲基礎，務求簡潔、自然。本方案以娛樂爲主。

\section{聲母}
普通聲母借用音值或系列相近的藏文字母。脣音、舌尖音、齒音和軟齶音分別使用藏文相應系列；小舌音借用腭音系列。抽象輔音 $C$、$N$、$S$ 只在文獻沒有進一步確定音值時使用。

\section{次要音節與冠音}
點號以前的次要音節另寫作一個短小字組。明寫元音照常標出；沒有明寫元音者加藏文 halanta，以別於明寫的 $ə$。緊密相連的前置輔音與主聲母組成同一字組。

\section{咽化、r 與合口}
咽化 $ˤ$、主聲母後的 $r$ 和合口 $ʷ$ 使用固定下加符號。$P$-$rV$ 與 $PrV$ 寫法相同；單獨的 $rV$ 仍以藏文 r 作基字。

\section{元音}
小寫 $a$ 不加符號；大寫 $A$ 明寫長元音符號。$i$、$u$、$e$、$o$、$ə$ 各用固定藏文元音符號。

\section{韻尾}
$j$、$w$、$r$、$m$、$n$、$ŋ$ 直接使用相應藏文字母。塞音韻尾 $p$、$t$、$k$ 依藏文後加字習慣寫作濁音字母。後加 $ʔ$、$s$、$ʔ$-$s$ 分別寫作藏文 a-chung、sa 及二者連寫。

\section{清響音}
Baxter--Sagart 2014 的 $m̥$、$n̥$、$ŋ̊$、$l̥$、$r̥$ 不另造藏文字母，而依其歷史來源寫成前置輔音加普通響音。

\section{拼寫測試}
代表性構形及其藏文寫法見網頁版。

\end{document}

